\documentclass[../main.tex]{subfiles}
\graphicspath{{\subfix{../images/}}}

\begin{document}
\section{Limitations}
This study has several limitations. The exposure assessment leveraged circuit-level data and we did not have access to individual-level exposure data. In the absence of more granular information about where individuals were on the circuits, to assign exposure we had to assume that individuals were evenly distributed across a circuit. This means that, particularly for larger circuits, we may have misclassified exposure if the distribution of people varied widely across the circuit. When we aggregated to the zip code, we may have misclassified exposed zip codes if individuals were not without power on the part of the circuit intersecting with a given zip code. Similarly, circuits may be partially energized and de-energized throughout a PSPS event. We did not have information on this, and therefore assumed that all of a circuit was out of power for the duration of a PSPS event. Again, this may have resulted in similar misclassification. 

Second, we relied on linking exposure data using residential zip code. As such, we could not ensure that an individual was (1) at home when a PSPS event occurred and (2) lived in the part of the zip code that was exposed to a PSPS event. 

Third, we are only capturing individuals who went to an emergency department or hospital. Individuals who seek healthcare may be systematically different than those who do not, and we are unable to determine the representativeness of our study population to all adults in California. 

Finally, in an observational study we cannot be certain that all confounding has been accounted for. Our case-crossover design captures much of the confounding, but longer term trends that we were unaware of may impact results.


\end{document}